\chapter{Conclusion}
\label{chap:conclusion}

A detailed profile from a later browser visit may explain that execution without identifying the cause of an earlier regression. The proposed research addresses this distinction by comparing follow-up profiling with diagnostic evidence retained during the original visit. Its central comparison concerns correct attribution of the original cause under matched resource allowances and a common deadline. The intended application is to help projects choose a performance profiler that balances monitoring cost, diagnostic accuracy and reliability within the tested conditions.

The five objectives connect repeatable measurement, a longitudinal dataset, independent fault labels, capture-strategy comparison and an assessment of reliability conditions. Persistent, transient and replaced-cause scenarios will distinguish recurrence from symptom similarity. Lightweight attribution will remain a strong baseline, and follow-up profiling will retain that original evidence. Bounded original capture and periodic capture will be evaluated alongside it, while always-rich capture will provide a separately costed reference. Missed alerts, failed captures, late answers, abstentions and policy exhaustion will remain in the verified eligible episode denominator.

The finite pilot will first establish whether fault checks, capture contracts, scheduling and complete resource accounting support the comparison. Its two fixture templates and repeated visits are a feasibility allocation, not a powered efficacy study or a representative sample of production websites. Any main experiment will require a separate protocol and analysis freeze, with pilot and repair data excluded from confirmatory evaluation. Existing work already establishes adaptive instrumentation, selective tracing and retention of original execution evidence, so the intended contribution is a bounded reliability and cost assessment rather than a claim that these mechanisms are new.

The work is currently at the literature review and research proposal stage. Neither the implementation, pilot nor main study has been completed. No diagnostic advantage, hypothesis result or final novelty claim can therefore be established from this report. The next stages are to finish the focused source review, implement and verify the measurement, fault and evidence contracts, and complete the execution manifest before the pilot. If feasibility is established, a separate main protocol will be frozen, followed by collection, analysis, discussion of limitations, identification of future work and the final empirical report. The resulting conclusions will be limited to the fixtures, faults, browser versions, states and resource conditions actually tested.
